\newpage
\section{Data Annotation Platform}
We use LabelStudio \cite{LabelStudio} as the platform for all our data annotations. We host LabelStudio in our internal clusters so that annotators can connect to the server conveniently via SSH-forwarding.
\subsection{Chart Selection}
\label{subsec:chart_selection_annotation}
\begin{figure}[h]
\centering
\includegraphics[width=0.75\linewidth]{figures/chart_annotation.jpg}
    \caption{Screenshot of our chart selection process. As shown in the screenshot, annotators are required to select one chart from 10 candidates figures that are pre-filtered with a cosine similarity $>0.65$ compared to the average chart embedding from MathVista.
    }
\vspace{-2ex}
\label{fig:chart_selection}
\end{figure}
\newpage

\subsection{Descriptive Question Annotation}
\label{subsec:descriptive_q_annotation}
\begin{figure}[h]
\centering
\includegraphics[width=1\linewidth]{figures/descriptive_annotation.png}
    \caption{Screenshot of our descriptive task annotation process. As shown in the screenshot, the annotator is presented with a chart and a randomly shuffled list of the 18 descriptive tasks (except Q19, which asks for the number of subplots and can be automatically converted from the number of subplot metadata) with GPT-generated answers. The annotator is required to select the first 3 answerable questions and the first unanswerable question with ground truth answers, fill in the number of subplots and the row, column number of the subplots to ask questions with (if the chart contains subplots).
    }
\vspace{-2ex}
\label{fig:descriptive_annotation}
\end{figure}
\newpage

\subsection{Reasoning Question Annotation}
\label{subsec:reasoning_q_annotation}
\begin{figure}[h]
\centering
\includegraphics[width=1\linewidth]{figures/reasoning_annotation.jpg}
    \caption{Screenshot of our reasoning task annotation process. As shown in the screenshot, the annotator is presented with a chart and a list of reasoning QAs automatically generated by GPT-4V. Then, the annotator needs to decide the final question to fill in (\textit{i.e.,} GPT-sourced, GPT-inspired, or human-written), and write down the final answer with an answer type (\textit{i.e.,} Text-in-Chart, Text-in-General, Number-in-Chart, Number-in-General). The answer type is subsequently used in the response generation process to provide additional instructions in generating response for the question.
    }
\vspace{-2ex}
\label{fig:reasoning_annotation}
\end{figure}
